\title{Group Sequence Policy Optimization}

\begin{abstract}
We present Group Sequence Policy Optimization (GSPO), a reinforcement learning algorithm specifically designed for robust, efficient, and effective training of large language models. In contrast to traditional methods that operate using token-level importance ratios, GSPO constructs its importance ratio using sequence-level likelihoods and implements clipping, reward calculation, and optimization at the sequence level. Our findings show that GSPO offers marked improvements in training speed and model performance relative to the GRPO algorithm, delivers enhanced stability during Mixture-of-Experts (MoE) RL training, and holds promise for reducing the complexity of RL system design. These strengths have played a critical role in the significant advancements realized in the latest Qwen3 models.
\end{abstract}

\section{Introduction}
\label{sec:intro}

Reinforcement learning (RL) has become a central approach for advancing the capabilities of language models at scale \citep{o1,dpsk-r1,qwq32b,qwen3}. Leveraging RL at large scale enables language models to address complex tasks—including advanced mathematics and programming—by fostering extended, more intricate reasoning sequences.

Achieving effective RL scaling with increased computational resources requires above all the preservation of steady and reliable training behavior. Yet, leading RL methods such as GRPO \citep{grpo} face profound challenges to stability when applied to enormous language models; this often culminates in dramatic and irreversible breakdowns of model performance \citep{qwen3, minimax-m1}. Such instability poses a significant obstacle to expanding the limits of what language models can accomplish through ongoing RL training.

In this study, we demonstrate that GRPO's instability originates from a fundamental flaw: the improper application and failure of importance sampling weights within its algorithm. This flaw results in substantial variance in training, generating noise that intensifies as response length grows and is exacerbated by the clipping process, eventually causing the model to collapse.

To mitigate these significant issues, we introduce \textbf{Group Sequence Policy Optimization (GSPO)}, a novel reinforcement learning algorithm designed for training large language models. GSPO's central contribution is its rigorously defined importance ratio, grounded in sequence likelihood as outlined by \citep{click}, which adheres to the essential principles of importance sampling. Moreover, GSPO utilizes normalized rewards calculated as advantages across multiple responses for each query, thereby reinforcing consistency between sequence-level reward signals and the optimization process.

Through our empirical analysis, we observe that GSPO substantially outperforms GRPO across measures of training stability, efficiency, and overall results. Importantly, GSPO intrinsically addresses the instability issues commonly encountered during reinforcement learning of large Mixture-of-Experts (MoE) models, thereby removing the necessity for elaborate stabilization mechanisms and suggesting possibilities for streamlining RL system architectures. These strengths of GSPO have directly facilitated notable performance gains in the most recent Qwen3 models. We anticipate that GSPO will serve as a reliable and scalable basis for further progress in large-scale RL training of language models.

\section{Preliminaries}

\paragraph{Notation}
Throughout this work, we represent an autoregressive language model with parameters $\theta$ as the policy $\pi_\theta$. The symbol $x$ is used for a query, while $\mathcal{D}$ refers to the set of all queries. For any given response $y$ to a query $x$, its probability under policy $\pi_\theta$ is expressed as $\pi_\theta (y | x)=\prod_{t=1}^{|y|} \pi_\theta (y_t | x, y_{<t} )$, where $|y|$ signifies the length of $y$ in tokens. Each query-response pair $(x, y)$ is evaluated by a verifier $r$, assigning a reward value $r(x, y) \in [ 0, 1 ]$.

\paragraph{Proximal Policy Optimization (PPO)} Using samples generated from the old policy $\pi_{\theta_\text{old}}$, PPO \citep{ppo} constrains the policy update within a proximal region of the old policy through the clipping mechanism. Specifically, PPO employs the following objective for policy optimization (we omit the KL regularization term hereinafter for brevity, as it is not the focus of this paper):

\begin{align}
\mathcal{J}_\text{PPO}(\theta) = \mathbb{E}_{ x \sim \mathcal{D},\, y \sim \pi_{\theta_\text{old}}( \cdot | x) }
\left[ \frac{1}{|y|} \sum_{t=1}^{|y|} 
\min \left( w_{t}(\theta) \widehat{A}_{t},  \, \mathrm{clip} \left( w_{t}(\theta), 1 - {\varepsilon}, 1 + {\varepsilon}\right) \widehat{A}_{t} \right)
\right],
\end{align}

Here, the token $y_t$'s importance ratio is expressed as $
w_{t}(\theta) = \frac{ \pi_{\theta} (y_{t} | x, y_{<t}) }{ \pi_{\theta_\text{old}} (y_{t} | x,y_{<t})} $, with the advantage $\widehat{A}_{t}$ for $y_t$ approximated by an additional value model, and $\varepsilon$ denoting the threshold used for clipping these importance ratios.

A central difficulty in practical applications of PPO arises from its significant dependence on the value model. Typically, this value model is constructed to match the policy model in scale, which leads to substantial demands on both memory and computational resources. Moreover, the algorithm's overall performance is fundamentally contingent upon the accuracy of the value estimates provided. Obtaining a trustworthy value model is itself a complicated task; extending its reliability to longer generations and more sophisticated scenarios further exacerbates the problem.

\paragraph{Group Relative Policy Optimization (GRPO)} GRPO \citep{grpo} bypasses the need for the value model by computing the relative advantage of each response within a group of responses to the same query. Specifically, GRPO optimizes the following objective:

\begin{align}
\label{equ:grpo}
\mathcal{J}_\text{GRPO}(\theta) = \mathbb{E}_{ x \sim \mathcal{D},\, \{y_i\}_{i=1}^G \sim \pi_{\theta_\text{old}}( \cdot | x) }
\left[ \frac{1}{G} \sum_{i=1}^{G} \frac{1}{|y_i|} \sum_{t=1}^{|y_i|} 
\min \left( w_{i,t}(\theta) \widehat{A}_{i,t},  \, \mathrm{clip} \left( w_{i,t}(\theta), 1 - {\varepsilon}, 1 + {\varepsilon}\right) \widehat{A}_{i,t} \right)
\right],
\end{align}

where $G$ is the number of generated responses to each query $x$ (i.e., the group size), and the importance ratio $w_{i,t}(\theta)$ and advantage $\widehat{A}_{i,t}$ of token $y_{i,t}$ are:

\begin{align}
    w_{i,t}(\theta)=\frac{ \pi_{\theta} (y_{i,t} | x, y_{i,<t}) }{ \pi_{\theta_\text{old}} (y_{i,t} | x,y_{i,<t})},\quad\ 
    \widehat{A}_{i,t} = \widehat{A}_{i} = \frac{r(x, y_i) - \mathrm{mean} \left( \{ r(x, y_i) \}_{i=1}^G \right) }{ \mathrm{std} \left( \{ r(x, y_i) \}_{i=1}^G \right) },
\end{align}

respectively, where all the tokens in $y_i$ share the same advantage as $\widehat{A}_{i}$.

\section{Motivation}
\label{sec:motivation}

The escalation in model parameters, incorporation of sparsity mechanisms (such as in Mixture-of-Experts architectures), and elongation of generated responses collectively demand the use of large rollout batch sizes to achieve efficient hardware throughput during reinforcement learning. In order to increase the efficiency of data usage, it is common to divide a substantial batch of rollout samples into several mini-batches for conducting multiple gradient update steps. However, this approach naturally gives rise to an off-policy learning context, as the responses $y$ used for updates are generated by an earlier version of the policy $\pi_{\theta_\text{old}}$ instead of the current, actively trained policy $\pi_\theta$. This situation underscores the importance of the clipping strategy employed in algorithms like PPO and GRPO, which serves to restrict the influence of samples that deviate excessively from the current policy during gradient calculations.

While mechanisms like clipping aim to manage this off-policy discrepancy, we identify a more fundamental issue in GRPO: \textit{its objective is ill-posed}. This problem becomes particularly acute when training large models on long-response tasks, leading to catastrophic model collapse. The ill-posed nature of the GRPO objective stems from a misapplication of importance sampling weights. The principle of importance sampling is to estimate the expectation of a function $f$ under a target distribution $\pi_\text{tar}$ by re-weighting samples drawn from a behavior distribution $\pi_\text{beh}$:

\begin{align}
\mathbb{E}_{z \sim \pi_\text{tar}} \left[ f(z) \right] = \mathbb{E}_{z \sim \pi_\text{beh}} \left[ \frac{ \pi_\text{tar} (z) }{ \pi_\text{beh} (z)} \, f(z) \right].
\end{align}

The effectiveness of the importance weight $\frac{ \pi_\text{tar} (z) }{ \pi_\text{beh} (z)}$ in correcting for discrepancies between distributions fundamentally depends on aggregating results from a substantial number of samples ($N \gg 1$) drawn from the behavioral distribution $\pi_\text{beh}$. This averaging process allows for proper adjustment to account for the distributional mismatch.

By contrast, GRPO introduces the importance weight $\frac{ \pi_{\theta} (y_{i,t} | x, y_{i,<t}) }{ \pi_{\theta_\text{old}} (y_{i,t} | x,y_{i,<t})}$ at every token $t$, using merely a single instance $y_{i,t}$ sampled from the corresponding next-token distribution $\pi_{\theta_\text{old}}( \cdot | x, y_{i, <t})$. This approach is insufficient for achieving meaningful distributional correction. Rather, it increases the variability in gradient estimates, with this noise intensifying as sequences become longer, particularly due to the clipping procedure. Our empirical investigations show this may precipitate catastrophic model collapse; such collapse proves difficult to reverse. Recovery efforts—such as reloading previous checkpoints, fine-tuning clipping thresholds, adjusting the length of generated sequences, or switching the RL querying strategy—are ultimately ineffective once collapse is observed.

These results highlight a pivotal flaw within the architecture of GRPO. The token-wise application of the importance weight overlooks a critical insight: \textbf{the granularity of the optimization target must correspond to the granularity of the reward function}. As rewards are administered to entire sequences, token-level off-policy adjustment is fundamentally ill-suited. This understanding drives us to abandon token-level correction, and instead, pursue sequence-level optimization and importance weighting that are directly aligned with the reward structure.

\section{Algorithm}

\subsection{GSPO: Group Sequence Policy Optimization}

Although the token-level importance weight $\frac{ \pi_{\theta} (y_{i,t} | x, y_{i,<t}) }{ \pi_{\theta_\text{old}} (y_{i,t} | x, y_{i,<t})}$ poses challenges within GRPO, we note that, for language generation tasks, the \textit{sequence-level} importance weight $\frac{ \pi_{\theta} (y | x) }{ \pi_{\theta_\text{old}} (y | x)}$ possesses a well-defined theoretical interpretation: it quantifies the divergence of a response $y$ drawn from $\pi_{\theta_\text{old}} (\cdot | x)$ from the distribution $\pi_{\theta} (\cdot | x)$. This property directly corresponds to the sequence-level reward and provides an informative metric for implementing the clipping procedure.

Based on this straightforward observation, we propose the \textbf{Group Sequence Policy Optimization (GSPO)} algorithm. GSPO employs the following sequence-level optimization objective:

\begin{align}
\mathcal{J}_\text{GSPO} (\theta) =
\mathbb{E}_{ x \sim \mathcal{D},\, \{y_i\}_{i=1}^G \sim \pi_{\theta_\text{old}}( \cdot | x) }
\left[ 
\frac{1}{G} \sum_{i=1}^{G}
\min \left( s_{i}(\theta)  \widehat{A}_{i},  \, \mathrm{clip} \left( s_{i}(\theta), 1 - {\varepsilon}, 1 + {\varepsilon} \right) \widehat{A}_{i} \right) 
\right],
\label{equ:gspo}
\end{align}

where we adopt the group-based advantage estimation:

\begin{align}
\widehat{A}_{i} = \frac{r(x, y_i) - \mathrm{mean} \left( \{ r(x, y_i) \}_{i=1}^G \right) }{ \mathrm{std} \left( \{ r(x, y_i) \}_{i=1}^G \right) },
\end{align}

and define the importance ratio $s_{i}(\theta)$ based on sequence likelihood \citep{click}:

\begin{align}
s_{i}(\theta) = \left( \frac{ \pi_{\theta} (y_i | x) }{ \pi_{\theta_\text{old}} (y_i | x)} \right)^{\frac{1}{|y_i|}}
=
\exp \left( \frac{1}{|y_i|} \sum_{t=1}^{|y_i|} \log \frac{ \pi_{\theta} (y_{i,t} | x, y_{i,<t}) }{ \pi_{\theta_\text{old}} (y_{i,t} | x,y_{i,<t})} \right).
\end{align}

Consequently, GSPO utilizes clipping on the whole response sequence rather than on individual tokens to filter out excessively ``off-policy'' samples during gradient calculation, aligning with both the reward and optimization processes at the sequence level. Additionally, we employ length normalization in $s_{i}(\theta)$ to mitigate variance and maintain $s_{i}(\theta)$ within a consistent numerical scope. Without such normalization, changes in the likelihoods of a small subset of tokens may induce significant instability in the sequence-level importance ratio, and responses of varying lengths would necessitate different clipping thresholds for their importance ratios. It is important to observe that GSPO's clipping bounds diverge substantially in scale from those used in prior approaches (such as GRPO), reflecting the differences in how importance ratios are formulated.

\subsection{Gradient Analysis}
\label{subsec:gradient}

We can derive the gradient of the GSPO objective as follows (clipping is omitted for brevity):

\begin{align}
\nabla_{\theta} \mathcal{J}_\text{GSPO} (\theta)
=&\ 
\nabla_{\theta} \mathbb{E}_{ x \sim \mathcal{D},\, \{y_i\}_{i=1}^G \sim \pi_{\theta_\text{old}}( \cdot | x) }
\left[ 
\frac{1}{G} \sum_{i=1}^{G} s_{i}(\theta) \widehat{A}_{i}
\right] \\
= &\ 
\mathbb{E}_{ x \sim \mathcal{D},\, \{y_i\}_{i=1}^G \sim \pi_{\theta_\text{old}}( \cdot | x) }
\left[ 
\frac{1}{G} \sum_{i=1}^{G}
s_{i}(\theta) \widehat{A}_{i}
\cdot \nabla_{\theta} \log s_{i}(\theta)
\right] \\
=&\ 
\mathbb{E}_{ x \sim \mathcal{D},\, \{y_i\}_{i=1}^G \sim \pi_{\theta_\text{old}}( \cdot | x) }
\left[ 
\frac{1}{G} \sum_{i=1}^{G} 
\left( \frac{ \pi_{\theta} (y_i | x) }{ \pi_{\theta_\text{old}} (y_i | x)} \right)^{\frac{1}{|y_i|}} \widehat{A}_{i} 
\cdot \frac{1}{|y_i|} \sum_{t=1}^{|y_i|} \nabla_{\theta} \log \pi_{\theta} (y_{i,t} | x, y_{i,<t}) 
\right].
\label{equ:gspo_grad}
\end{align}

For comparison, the gradient of the GRPO objective is as follows (note that $\widehat{A}_{i,t} = \widehat{A}_{i}$):

\begin{align}
\nabla_{\theta} \mathcal{J}_\text{GRPO} (\theta) 
=&\ 
\nabla_{\theta} \mathbb{E}_{ x \sim \mathcal{D},\, \{y_i\}_{i=1}^G \sim \pi_{\theta_\text{old}}( \cdot | x) }
\left[ \frac{1}{G} \sum_{i=1}^{G} \frac{1}{|y_i|} \sum_{t=1}^{|y_i|} 
w_{i,t}(\theta) \widehat{A}_{i,t}
\right] \\
=&\
\mathbb{E}_{ x \sim \mathcal{D},\, \{y_i\}_{i=1}^G \sim \pi_{\theta_\text{old}}( \cdot | x) }
\left[ \frac{1}{G} \sum_{i=1}^{G} \widehat{A}_{i} 
\cdot \frac{1}{|y_i|} \sum_{t=1}^{|y_i|} 
\frac{ \pi_{\theta} (y_{i,t} | x, y_{i,<t}) }{ \pi_{\theta_\text{old}} (y_{i,t} | x,y_{i,<t})} 
\nabla_{\theta} \log \pi_{\theta} (y_{i,t} | x, y_{i,<t})  
\right].
\end{align}

Consequently, the core difference between GSPO and GRPO is rooted in \textit{the manner in which they assign weights to the gradients associated with token log likelihoods}. GRPO applies individualized ``importance weights'' to tokens, expressed as $\frac{ \pi_{\theta} (y_{i,t} | x, y_{i,<t}) }{ \pi_{\theta_\text{old}} (y_{i,t} | x,y_{i,<t})}$. These varying weights—spanning the intervals $(0, 1+\varepsilon]$ when $\widehat{A}_{i} >0$ and $[1-\varepsilon, +\infty)$ when $\widehat{A}_{i} < 0$—have significant, non-trivial effects that can aggregate throughout training, potentially producing unforeseen outcomes. By contrast, GSPO uniformly weights all tokens within a response, thereby removing the source of instability inherent in the GRPO approach.

\subsection{GSPO-token: A Token-level Objective Variant}

In scenarios like multi-turn RL, we may desire a finer-grained advantage adjustment than the sequence level. To this end, we introduce a token-level objective variant of GSPO, namely \textbf{GSPO-token}, to allow token-wise advantage customization:

\begin{align}
\mathcal{J}_\text{GSPO-token}(\theta)
=
\mathbb{E}_{ x \sim \mathcal{D},\, \{y_i\}_{i=1}^G \sim \pi_{\theta_\text{old}}( \cdot | x) }
\left[ 
\frac{1}{G} \sum_{i=1}^{G} \frac{1}{|y_i|} \sum_{t=1}^{|y_i|} 
\min \left( s_{i,t}(\theta) \widehat{A}_{i,t},  \, \mathrm{clip} \left( s_{i,t}(\theta), 1 - {\varepsilon}, 1 + {\varepsilon} \right) \widehat{A}_{i,t} \right) 
\right],
\label{equ:gspo-token}
\end{align}

where

\begin{align}
s_{i,t}(\theta) 
= \mathrm{sg} \left[ s_{i}(\theta) \right]  \cdot \frac{ \pi_{\theta} (y_{i,t} | x, y_{i,<t}) }{ \mathrm{sg} \left[ \pi_{\theta} (y_{i,t} | x, y_{i,<t}) \right] },
\end{align}

and $\mathrm{sg}[\cdot]$ denotes only taking the numerical value but stopping the gradient, corresponding to the \texttt{detach} operation in PyTorch. The gradient of GSPO-token can be derived as:

\begin{align}
\nabla_{\theta} \mathcal{J}_\text{GSPO-token}(\theta)
=&\ 
\nabla_{\theta} \mathbb{E}_{ x \sim \mathcal{D},\, \{y_i\}_{i=1}^G \sim \pi_{\theta_\text{old}}( \cdot | x) }
\left[ 
\frac{1}{G} \sum_{i=1}^{G}
\frac{1}{|y_i|} \sum_{t=1}^{|y_i|} 
s_{i,t}(\theta) \widehat{A}_{i,t}
\right] \\
=&\ 
\mathbb{E}_{ x \sim \mathcal{D},\, \{y_i\}_{i=1}^G \sim \pi_{\theta_\text{old}}( \cdot | x) }
\left[ 
\frac{1}{G} \sum_{i=1}^{G} s_i (\theta) \cdot \frac{1}{|y_i|} \sum_{t=1}^{|y_i|} 
\widehat{A}_{i,t} \frac{ \nabla_{\theta} \pi_{\theta} (y_{i,t} | x, y_{i,<t}) }{ \pi_{\theta} (y_{i,t} | x, y_{i,<t}) }
\right] \\
=&\ 
\mathbb{E}_{ x \sim \mathcal{D},\, \{y_i\}_{i=1}^G \sim \pi_{\theta_\text{old}}( \cdot | x) }
\left[ 
\frac{1}{G} \sum_{i=1}^{G}  \left( \frac{ \pi_{\theta} (y_i | x) }{ \pi_{\theta_\text{old}} (y_i | x)} \right)^{\frac{1}{|y_i|}}
\cdot \frac{1}{|y_i|} \sum_{t=1}^{|y_i|}
\widehat{A}_{i,t} \nabla_{\theta} \log \pi_{\theta} (y_{i,t} | x, y_{i,<t})  
\right].
\label{equ:gspo-token_grad}
\end{align}

It should be observed that $\frac{ \pi_{\theta} (y_{i,t} | x, y_{i,<t}) }{ \mathrm{sg} \left[ \pi_{\theta} (y_{i,t} | x, y_{i,<t}) \right] }$ evaluates to 1. As a consequence, $s_{i,t}(\theta)$ and $s_{i}(\theta)$ produce the same numerical result. By examining Equation~\eqref{equ:gspo} alongside Equation~\eqref{equ:gspo-token}, and comparing Equation~\eqref{equ:gspo_grad} with Equation~\eqref{equ:gspo-token_grad}, one can see that both GSPO-token and GSPO yield equivalent values for the objective function, clipping criterion, and the formal gradient, provided the advantage values are set uniformly for every token in $y_i$ (specifically, $\widehat{A}_{i,t} = \widehat{A}_{i}$). Nevertheless, GSPO-token has the added benefit of permitting individual advantage assignments for each token, granting greater adaptability.

\section{Experiments and Discussion}

\subsection{Empirical Results}

We conduct experiments utilizing a cold-start model initialized from Qwen3-30B-A3B-Base, presenting both the training reward trajectories and performance trends on various benchmarks: AIME'24 (average Pass@1 across 32 samples), LiveCodeBench (202410-202502, average Pass@1 over 8 samples), and CodeForces (Elo Rating). Throughout the RL training process, each batch of collected rollout data is split into four mini-batches to facilitate gradient update steps. For GSPO, we configure the clipping intervals in Equation~\eqref{equ:gspo} to 3e-4 (left) and 4e-4 (right). In comparison, GRPO is employed as the baseline, where the clipping thresholds in Equation~\eqref{equ:grpo} are set to 0.2 (left) and 0.27 (right)—values carefully selected for equitable evaluation. It is important to highlight that GRPO requires the use of the Routing Replay training mechanism to achieve stable convergence in MoE RL, which we elaborate on further in \S~\ref{subsec:routing_replay}. In contrast, \textit{GSPO removes the necessity for this additional strategy}.

As illustrated in Figure~\ref{fig:results}, GSPO training maintains consistent stability across all stages. Our findings indicate that \textbf{GSPO facilitates ongoing performance gains by scaling training compute, periodically refreshing the query set, and lengthening the output generation}. In addition, GSPO exhibits enhanced training efficiency compared to GRPO, attaining higher training accuracy and benchmark outcomes with equivalent training compute and query consumption. Ultimately, GSPO has been effectively implemented in RL training for the recent Qwen3 models, providing strong evidence of GSPO's effectiveness in harnessing RL scaling for large language models.

\subsection{Curious Observation on Clipping Fractions}

An important difference between GSPO and GRPO lies in GSPO's approach to clipping, which operates at the level of entire responses rather than on individual tokens. As illustrated in Figure~\ref{fig:clipping}, the proportion of tokens clipped under GSPO is two orders of magnitude greater compared to GRPO, and this substantial gap remains unaffected by modifications to the clipping ranges. Interestingly, although GSPO discards a much larger number of tokens and subsequently uses a smaller dataset for model training or gradient estimation, it nonetheless attains superior training efficiency relative to GRPO. This unexpected outcome — where aggressive token clipping translates into improved training efficiency — highlights the intrinsic noise and ineffectiveness in GRPO’s token-level gradient estimation for sample utilization. Conversely, GSPO's sequence-oriented method generates a more consistent and beneficial signal for learning.

\subsection{Benefit of GSPO for MoE Training}
\label{subsec:routing_replay}

\paragraph{Background}
Unlike RL training in dense models, the inherently sparse activation patterns of MoE models pose distinct challenges for maintaining training stability. Notably, we observed that when applying the GRPO algorithm, fluctuations in \textit{expert activation} across updates — referred to as \textit{expert-activation volatility} — undermine the convergence of RL training in MoE architectures. Specifically, following one or more steps of gradient optimization, the set of experts engaged for identical responses can vary substantially. For instance, for the 48-layer Qwen3-30B-A3B-Base architecture, analysis after each RL-driven update indicates that approximately 10\% of experts chosen under the revised policy $\pi_{\theta}$ differ from those selected by the previous policy $\pi_{\theta_\text{old}}$ for a given rollout instance. This issue, which intensifies with increased model depth in MoEs, leads to considerable oscillations in the token-level importance ratios $w_{i,t}(\theta) = \frac{ \pi_{\theta} (y_{i,t} | x, y_{i,<t}) }{ \pi_{\theta_\text{old}} (y_{i,t} | x,y_{i,<t})}$, thereby diminishing their reliability (as elaborated in \S~\ref{sec:motivation} and~\ref{subsec:gradient}), and ultimately disrupting the expected RL training convergence behavior.

\paragraph{Our Previous Approach} In our earlier work, we addressed this issue by utilizing the \textbf{Routing Replay} training method. More concretely, during training, we store the set of experts activated under $\pi_{\theta_\text{old}}$, and subsequently ``replay'' these activation patterns with $\pi_{\theta}$ for the calculation of the importance ratios $w_{i,t}(\theta) = \frac{ \pi_{\theta} (y_{i,t} | x, y_{i,<t}) }{ \pi_{\theta_\text{old}} (y_{i,t} | x,y_{i,<t})}$. By doing so, for every token $y_{i,t}$, both $\pi_{\theta} (y_{i,t} | x, y_{i,<t})$ and $\pi_{\theta_\text{old}} (y_{i,t} | x,y_{i,<t})$ utilize identical expert configurations, thereby maintaining stable token-level importance ratios and allowing for the consistent optimization of the activated network throughout successive gradient steps. As shown in Figure~\ref{fig:routing_replay}, Routing Replay plays a pivotal role in ensuring proper convergence during the GRPO training of MoE models.

\paragraph{Benefit of GSPO} While Routing Replay facilitates successful GRPO training in MoE models, it introduces extra memory usage and communication costs due to its reliance on reused routing decisions, which in turn restricts the MoE model's effective capacity. In comparison, as illustrated in Figure~\ref{fig:results}, GSPO discards the necessity for Routing Replay, readily calculates importance ratios $s_i(\theta)$ in the standard way, achieves steady convergence, and ensures stable optimization. The central observation is that GSPO prioritizes sequence-level likelihoods (specifically, $\pi_{\theta} (y_i | x)$) rather than focusing on token-level likelihoods (such as $\pi_{\theta} (y_{i,t} | x, y_{i,<t})$), making it robust to variations at the token level. Because MoE models consistently retain strong language modeling abilities, their sequence likelihoods exhibit little fluctuation. Collectively, GSPO directly addresses and mitigates the expert activation instability in MoE models, thus removing the necessity for intricate solutions like Routing Replay. As a consequence, GSPO enhances training simplicity and reliability, while enabling the model to fully utilize its designed capacity without imposed limitations.

\subsection{Benefit of GSPO for RL Infrastructure}

Due to precision mismatches that may arise between training engines such as Megatron and inference systems like SGLang or vLLM, it is standard practice to recalculate sampled response likelihoods under the old policy $\pi_{\theta_\text{old}}$ using the training engine. Nevertheless, since GSPO employs sequence-level likelihoods instead of token-level ones for its optimization procedure, it is generally less affected by such discrepancies. As a result, GSPO enables the optimization process to make direct use of likelihoods provided by the inference engine, eliminating the need for recalculation with the training engine. This direct approach is particularly advantageous in cases such as partial rollout, multi-turn reinforcement learning, and within training-inference decoupled architectures.

\section{Conclusion}

Group Sequence Policy Optimization (GSPO) is introduced as a novel reinforcement learning method specifically tailored for optimizing large language models. GSPO applies the principles of importance sampling, calculating importance ratios via sequence likelihoods and executing sequence-level operations such as clipping, rewarding, and optimization. The algorithm yields remarkable advantages in training robustness, efficiency, and overall outcomes over GRPO, proving particularly effective in large-scale RL applications with MoE models and underpinning the significant progress observed in the newest Qwen3 models. As GSPO represents a scalable foundation for algorithmic development, ongoing efforts will focus on further expanding RL capabilities, anticipating substantial progress in artificial intelligence.

\end{document}